% -------------------------------------------------
% 3) DIAPOSITIVA DE CONTENIDO (logo abajo derecha)
% -------------------------------------------------
\begin{frame}{Título de la diapositiva de contenido}
  \begin{block}{Idea principal / Resultado}
    Texto placeholder para explicar la idea principal. Podés escribir 2--4 líneas
    sin que quede demasiado aire.
  \end{block}

  \vspace{0.25cm}

  \textbf{Puntos clave:}
  \begin{itemize}
    \item Punto 1 (placeholder).
    \item Punto 2 (placeholder).
    \item Punto 3 (placeholder).
    \item hola
  \end{itemize}

  \vspace{0.15cm}

  \textbf{Figura (placeholder):}
  \begin{center}
    \fbox{
      \parbox{0.75\linewidth}
      {
        \centering 
          \begin{figure}
            \includegraphics[width=0.8\textwidth]{./attatchments/unc3_a.jpg}
            \caption{Figura placeholder. Reemplazá esto con tu propia figura.}
          \end{figure}
      }
      }
  \end{center}
\end{frame}

\begin{frame}{Título de la diapositiva con imagen}
  \begin{columns}[T,onlytextwidth]

    % ---- Columna izquierda: texto ----
    \column{0.45\textwidth}
    \textbf{Idea principal}
    Texto
    % ---- Columna derecha: imagen ----
    \column{0.55\textwidth}
    \begin{center}
      \includegraphics[
        width=\linewidth,
        height=0.75\textheight,
        keepaspectratio
      ]{./attatchments/unc3_a.jpg}
    \end{center}

  \end{columns}
\end{frame}

\begin{frame}{Solo imagen}
  \centering
    \includegraphics[
      width=\linewidth,
      height=0.82\textheight,
      keepaspectratio
    ]{./attatchments/imagen.png}
\end{frame}